\documentclass[12pt, a4paper]{ctexart}
\usepackage{geometry}
\geometry{left=2.5cm, right=2.5cm, top=3cm, bottom=3cm}
\usepackage{listings}
\usepackage{graphicx}
\usepackage{xcolor}
\usepackage{amsmath}
\usepackage{booktabs}
\usepackage{float}
\usepackage{tikz}
\usetikzlibrary{positioning, arrows.meta}

% Code snippet settings
\lstset{
    basicstyle=\ttfamily\small,
    frame=single,
    framesep=5pt,
    xleftmargin=10pt,
    xrightmargin=10pt,
    keywordstyle=\color{blue},
    commentstyle=\color{gray},
    breaklines=true
}

\begin{document}

% 封面
\begin{titlepage}
    \centering
    \vspace*{5cm}
    {\Huge\bfseries 实验报告\par}
    \vspace{2cm}
    {\Large 课程：数字逻辑与计算机组成实验\par}
    \vspace{1cm}
    {\Large 实验十二：计算机系统\par}
    \vspace{4cm}
    {\large 姓名：\underline{\makebox[4cm]{郑袭明}}\par}
    \vspace{1cm}
    {\large 学号：\underline{\makebox[4cm]{251502027}}\par}
    \vspace{1cm}
    {\large 日期：\today\par}
\end{titlepage}

\tableofcontents
\newpage

\section{系统概述}

这次实验要做的是一台能真正用起来的小计算机。它的核心是前面几次实验里搭好的那个
RV32I 五级流水线 CPU，本身已经支持完整的整数指令、数据前递、load-use 停顿，以及
按字节、半字、字访存。但只有 CPU 还不够——它得能把东西显示出来，能接受输入，还得
有个时间概念。于是这次给它接上了三样外设：VGA 字符显示、PS/2 键盘和一个秒级定时器，
再在上面写一套裸机软件，让人可以坐在屏幕前敲命令。

\begin{table}[h]
\centering
\begin{tabular}{ll}
\toprule
文件 & 作用 \\
\midrule
\lstinline|top.v| & 顶层，连接 CPU、存储器与各外设 \\
\lstinline|rv32ip.v| & RV32I 五级流水线 CPU \\
\lstinline|alu.v|、\lstinline|regfile.v| & ALU 与寄存器堆 \\
\lstinline|dmem.v| & 数据存储器，含字节/半字访存 \\
\lstinline|vga_char.v|、\lstinline|vga_ctrl.v| & VGA 字符显存与扫描时序 \\
\lstinline|keyboard.v| 等 & PS/2 键盘解析与扫描码转换 \\
\lstinline|timer.v| & 秒级定时器 \\
\bottomrule
\end{tabular}
\end{table}

很多模块之前已经做完了，这里真正需要做的，是把这些外设挂到 CPU
的数据总线上，也就是 \lstinline|top.v| 里那套内存映射 I/O 的接入逻辑。软件部分则是
从头写的：一个带管道和重定向的 shell、一个内存文件系统、行编辑与输出系统、显示和键盘
的底层驱动。

\section{硬件：把外设接到 CPU 数据总线上}

\subsection{内存映射 I/O 的总体接入}

CPU 对外只有一组数据存储器端口：地址 \lstinline|dmemaddr|、写数据 \lstinline|dmemdatain|、
读数据 \lstinline|dmemdataout|、写使能 \lstinline|dmemwe|、访存类型 \lstinline|dmemop|。
读在时钟上升沿、写在下降沿。问题在于，现在挂在这组端口后面的不止一块 RAM，还有三个外设。
所以现在将一组总线信号按地址分段。用地址的高 12 位 \lstinline|dmemaddr[31:20]| 当作设备选择，
不同的段对应不同的设备：

\begin{table}[H]
\centering
\begin{tabular}{llll}
\toprule
段（\lstinline|[31:20]|） & 基址 & 设备 & 读/写 \\
\midrule
\lstinline|0x001| & \lstinline|0x00100000| & 数据 RAM & 读写 \\
\lstinline|0x002| & \lstinline|0x00200000| & VGA 字符缓冲 & 写 \\
\lstinline|0x003| & \lstinline|0x00300000| & 键盘寄存器 & 读 \\
\lstinline|0x004| & \lstinline|0x00400000| & 定时器 & 读写 \\
\bottomrule
\end{tabular}
\end{table}

写一侧的处理是把 \lstinline|dmemwe| 按段拆开，每个设备各自得到一个只在本段地址下才有效的写使能，
这样一次写绝不会同时落到两个设备上：

\begin{lstlisting}[language=Verilog]
wire dmem_we = dmemwe & (dmemaddr[31:20] == 12'h001);  // 数据 RAM
wire vga_we  = dmemwe & (dmemaddr[31:20] == 12'h002);  // VGA
wire time_we = dmemwe & (dmemaddr[31:20] == 12'h004);  // 定时器
\end{lstlisting}

读一侧则反过来，用一个多路选择器按地址段把对应设备的输出汇成 \lstinline|dmemdataout|
送回 CPU：

\begin{lstlisting}[language=Verilog]
assign dmemdataout = (dmemaddr[31:20] == 12'h001) ? dmem_out :
                     (dmemaddr[31:20] == 12'h003) ? key_reg  :
                     (dmemaddr[31:20] == 12'h004) ? time_reg :
                                                    32'b0    ;
\end{lstlisting}

剩下的就是在 \lstinline|top.v| 里例化各外设，把地址、数据、时钟这些信号接到它们的端口上。

\subsection{接入 VGA 显示}

VGA 这块直接例化了两个模块：\lstinline|vga_char| 管显存和字符渲染，\lstinline|vga_ctrl|
管扫描时序和 RGB 输出。接入时，把上一节算出的 \lstinline|vga_we|、地址 \lstinline|dmemaddr[16:0]|
和写数据的低 8 位接到 \lstinline|vga_char| 的写口；\lstinline|vga_char| 输出的像素数据再接给
\lstinline|vga_ctrl|，由它驱动板上的 VGA 引脚。像素时钟是 25\,MHz，用一个时钟向导 IP
\lstinline|clk_wiz_0| 从 100\,MHz 主时钟分频得到。

\begin{lstlisting}[language=Verilog]
wire vga_we = dmemwe & (dmemaddr[31:20] == 12'h002);
vga_char vga (
    .wclk (dmemwrclk), .we (vga_we),
    .waddr(dmemaddr[16:0]), .wdata(dmemdatain[7:0]),
    .h_addr(h_addr), .v_addr(v_addr), .vga_data(vga_data)
);
\end{lstlisting}

每个字符单元的线性地址是“行左移 7 位再加列”，也就是行间距 128（实际只用前 80 列）。
除了写字符，显存还藏着两个控制寄存器，靠写地址的高位来区分写的是哪一种：

\begin{lstlisting}[language=Verilog]
if (waddr[16]) start_row  <= wdata[4:0]; // 滚动寄存器：屏幕起始行
else if (waddr[15]) cursor_pos <= waddr[11:0]; // 光标位置寄存器
else vram[waddr[11:0]] <= wdata; // 普通写字符
\end{lstlisting}

也就是说，写 \lstinline|0x00210000| 一带改的是滚动起点（用来整屏上滚），写
\lstinline|0x00208000| 一带改的是光标位置，其余地址才是往显存里填字符。

\subsection{接入 PS/2 键盘}

键盘的 \lstinline|keyboard| 模块有两个需要的输出：当前按键的 ASCII 码 \lstinline|ascii_key|，
以及一个每按一次键就加一的计数 \lstinline|key_count|。接入要做的是把这两个值拼成一个
32 位字，在读时钟沿锁存进 \lstinline|key_reg|，CPU 读 \lstinline|0x003| 段时拿到的就是它：

\begin{lstlisting}[language=Verilog]
keyboard kbd (
    .clk(CLK100MHZ), .clrn(CPU_RESETN),
    .ps2_clk(PS2_CLK), .ps2_data(PS2_DATA),
    .key_count(key_count), .ascii_key(ascii_key)
);
reg [31:0] key_reg;
always @(posedge dmemrdclk) key_reg <= {16'b0, key_count, ascii_key};
\end{lstlisting}

这里 \lstinline|key_count| 是关键。键盘没有缓冲队列，软件没法直接问“有没有新键”，于是
用这个计数来判断：软件记住上一次读到的计数，发现计数变了就说明又按了一次键，这时
\lstinline|ascii_key| 里的就是新字符。后面 \lstinline|getch| 正是这么做的。另外，方向键、
Home、End 这类键被映射成了约定的控制码（比如左方向对应 \lstinline|0x02|），软件
\lstinline|io.c| 里的按键定义就是按这套来的。

\subsection{接入定时器}

定时器最简单。例化已有的 \lstinline|timer| 模块，把按段拆出来的写使能 \lstinline|time_we|
和写数据接上去，它的输出 \lstinline|time_reg| 已经并进了上面那个读多路里。\lstinline|time_reg|
是一个以秒为单位的计数：模块内部在 100\,MHz 下分频成 1\,Hz 自增。软件往 \lstinline|0x004|
段写一个值就能直接给它置数，也就是“对时”，这正是软件里 \lstinline|time set HH:MM:SS|
的硬件依据。

\section{软件：裸机系统与 Shell}

软件这部分是这次从头写的，分几层：最底下是显示和键盘的驱动（\lstinline|sys.c|），往上是
行编辑和输出系统（\lstinline|io.c|），再往上有字符串工具（\lstinline|str.c|）和一个内存
文件系统（\lstinline|fs.c|），最上层是把这些拼起来的 shell（\lstinline|main.c|）。整个程序
是裸机运行的，没有操作系统也没有 C 标准库，所以进 \lstinline|main| 之前还得自己用一小段
汇编把栈指针摆好（\lstinline|main.c:9-13|）。下面按从底到上的顺序讲。

\subsection{显示与键盘底层驱动（\texttt{sys.c}）}

驱动要做的就是把上一章那些硬件约定翻译成 C 里能调的函数。先把显存的三个地址用指针表示
出来，对应显存、光标寄存器、滚动寄存器（\lstinline|sys.c:3-5|）：

\begin{lstlisting}[language=C]
char *vga_start = (char *)VGA_START;       // 0x00200000 字符缓冲
volatile char *vga_cursor = (volatile char *)VGA_CURSOR;  // 光标
volatile int  *vga_lineo  = (volatile int *)VGA_LINE_O;   // 滚动
\end{lstlisting}

往屏幕打一个字符由 \lstinline|putch| 完成（\lstinline|sys.c:45-67|）。它要分几种情况：
回车换行就调 \lstinline|new_line|；退格就把光标往回退一格并擦掉那个字符；
普通字符就写进显存当前位置，再把列号加一，写满一行也自动换行。每动一次都要更新一下
硬件光标的位置。

滚屏是这里比较有意思的地方。屏幕只有 30 行，写满之后不能停，得让内容往上滚。我的做法不是
真的把显存整体搬一遍，而是改硬件的滚动寄存器：维护一个起始行 \lstinline|start_line|，
写满时让它往前挪一行并清掉新腾出来的那行，然后把新的起始行写进 \lstinline|*vga_lineo|
（\lstinline|sys.c:16-28|）。硬件按这个起始行去取显存，看起来就是整屏向上滚了一行，而
搬运的活完全交给了硬件。

读键盘是 \lstinline|getch|（\lstinline|sys.c:69-80|）。它就是不停地读键盘寄存器：

\begin{lstlisting}[language=C]
char getch(void) {
    static int last_count = 0;
    volatile unsigned int *key_reg = (volatile unsigned int *)KEY_START;
    while (1) {
        unsigned int w = *key_reg;
        int count = (int)((w >> 8) & 0xff);
        if (count != last_count) {   // 计数变了，说明有新键
            last_count = count;
            return (char)(w & 0xff);
        }
    }
}
\end{lstlisting}

计数没变就一直转，一旦发现计数和上次不一样，就取出低 8 位的 ASCII 返回。

\subsection{行编辑与输出系统（\texttt{io.c}）}

\lstinline|getch| 只能拿到一个字符，但用户敲命令时希望能改、能移动光标。\lstinline|read_line|
（\lstinline|io.c:20-78|）就是为此写的一个完整行编辑器：支持左右移动、Home 跳到行首、
End 跳到行尾、退格、Delete，以及在行中间插入字符。它维护两个量——缓冲里的内容长度
\lstinline|len|，和光标当前所在的位置 \lstinline|pos|。插入或删除字符时，要把后面的内容在
缓冲里整体挪一格，再用 \lstinline|redraw| 把光标之后的部分重画一遍，保证屏幕和缓冲始终
一致。

这里还做了一件对 shell 很关键的事：输出重定向。平时 \lstinline|print| 直接往屏幕写，但
有时需要把一条命令的输出先攒到内存里（比如要喂给管道的下一段，或者要写进文件）。于是
设了一个“输出去向”的开关（\lstinline|io.c:81-108|）：

\begin{lstlisting}[language=C]
void out_redirect(char *buf, int cap) { /* 后续输出改写到 buf */ }
void out_restore(void) { /* 恢复成直接写屏幕 */ }
static void out_char(char c) {
    if (out_buf) { /* 写进缓冲并补 '\0' */ }
    else putch(c); // 否则照常上屏
}
\end{lstlisting}

所有打印最终都走 \lstinline|out_char|，所以只要切一下这个开关，同一套打印代码既能上屏
也能写进缓冲。

打印接口本身也做得顺手一些。\lstinline|io.h| 里用 C11 的 \lstinline|_Generic| 按参数类型
自动挑函数：字符串走 \lstinline|print_str|，整数走 \lstinline|print_int|，无符号走
\lstinline|print_uint|，字符走 \lstinline|print_char|；再配一组按参数个数展开的宏，支持一次
传 1 到 10 个参数。于是上层可以直接写 \lstinline|println("x=", n)| 这种混合类型的打印，
算是一个不实现 \lstinline|printf| 能做到的简化版本。

\subsection{字符串工具（\texttt{str.c}）}

\lstinline|str.c| 是一组小工具，没有标准库只能自己写：\lstinline|str_eq| 比较两个字符串是否
相等，\lstinline|str_to_i| 把数字串转成整数，\lstinline|str_len| 求长度，\lstinline|str_copy|
做限长拷贝（防越界），\lstinline|str_contains| 在一个串里找子串。它们被上层到处用到，
比如命令分发要比较命令名，\lstinline|grep| 要找子串。

\subsection{内存文件系统（\texttt{fs.c}）}

文件系统是纯内存的，所有东西都放在 RAM 里，掉电就没了。它用一个固定 64 个节点的数组当
存储，每个节点要么是目录要么是文件，靠三个下标 \lstinline|parent|、\lstinline|first_child|、
\lstinline|next_sibling| 串成一棵目录树，文件内容就存在节点里那块 256 字节的
\lstinline|data| 里（\lstinline|fs.c:9-22|）：

\begin{lstlisting}[language=C]
typedef struct {
    char name[16];
    unsigned char used, is_dir;
    int parent, first_child, next_sibling;  // 用下标串成树
    const char *content;
    char data[256];                          // 文件内容内联存放
} fs_node;
\end{lstlisting}

用“孩子—兄弟”表示法，每个目录只记第一个孩子，孩子之间用兄弟指针连成一条链，这样一棵
任意宽度的树只用三个下标就能表达。

路径解析是 \lstinline|fs_resolve|（\lstinline|fs.c:73-98|）。它从根目录（路径以 \lstinline|/|
开头）或当前目录出发，按 \lstinline|/| 把路径切成一段段，逐段在当前目录的孩子里找名字
匹配的节点；遇到 \lstinline|.| 原地不动，遇到 \lstinline|..| 回到父目录。创建文件或目录时
还要先用 \lstinline|split_parent|（\lstinline|fs.c:162-182|）把路径拆成“父目录”加“最后那一
段名字”。

在这之上就是各个命令：\lstinline|pwd| 顺着父链往上走再倒着打出来；\lstinline|cd| 换当前
目录；\lstinline|ls| 列出一个目录下的孩子（目录名后加 \lstinline|/|）；\lstinline|cat| 打文件
内容；\lstinline|mkdir| 建目录；\lstinline|rm| 删文件；\lstinline|rmdir| 删空目录，并且会挡住
删根目录、删当前目录、删非空目录这几种情况。给 shell 用的两个接口是 \lstinline|fs_read|
和 \lstinline|fs_write|（\lstinline|fs.c:238-271|）：读返回内容指针；写支持新建、覆盖，也
支持追加（\lstinline|append|），分别对应重定向的 \lstinline|>| 和 \lstinline|>>|。

\subsection{Shell（\texttt{main.c}）}

最上层把前面这些拼成一个能交互的 shell。主循环很短：打印提示符 \lstinline|$ |，读一行，
执行，循环（\lstinline|main.c:283-293|）。

所有命令集中登记在一张表里，用的是 X 宏的写法（\lstinline|main.c:123-163|）。命令名和处理
函数只在 \lstinline|COMMAND_LIST| 里写一遍，再分别展开成“注册”和“分发”两处代码，以后加
命令只改这一个地方：

\begin{lstlisting}[language=C]
#define COMMAND_LIST(X) \
    X("hello", cmd_hello) X("clear", cmd_clear) \
    X("time", cmd_time)   X("echo", cmd_echo)   \ 
    X("grep", cmd_grep)   X("ls", cmd_ls) // ...
// 分发时再展开一次
#define X(name, fn) if (str_eq(cmd, name)) { fn(args); return; }
    COMMAND_LIST(X)
#undef X
\end{lstlisting}

一行命令真正的处理在 \lstinline|run_command|（\lstinline|main.c:230-281|），它做了三件事。
先按竖线 \lstinline@|@ 把整行切成若干段（最多 8 段），这就是管道；每一段先去掉首尾空格，再用
\lstinline|parse_redir|（\lstinline|main.c:176-225|）把 \lstinline|<|、\lstinline|>|、\lstinline|>>|
这些重定向解析出来；然后顺序执行每一段。段与段之间的数据传递用两块缓冲 \lstinline|bufA|
和 \lstinline|bufB| 轮流当中转：前一段的输出存进缓冲，作为后一段的输入 \lstinline|g_stdin|，
最后一段则直接输出到屏幕。如果某段带了 \lstinline|>| 或 \lstinline|>>|，就先把输出攒进
\lstinline|filebuf|，执行完再写进文件系统。这套流程靠的正是 \lstinline|io.c| 那个输出重定向
开关。

内建命令本身都不复杂：\lstinline|hello| 和 \lstinline|clear| 一目了然；\lstinline|time| 读出
当前时间，\lstinline|time set HH:MM:SS| 则用 \lstinline|parse_hms| 解析后写定时器对时；
\lstinline|fibn| 迭代算第 n 个斐波那契数；\lstinline|echo| 原样回显；\lstinline|grep| 在输入
\lstinline|g_stdin| 上逐行匹配，配合管道用。串起来就能敲出这样的命令行：

\begin{lstlisting}
$ echo hello | grep he
hello
$ ls > a.txt
$ cat a.txt
\end{lstlisting}

第一条把 \lstinline|echo| 的输出通过管道喂给 \lstinline|grep|；第二条把 \lstinline|ls| 的结果
重定向写进文件；第三条再读出来。管道、重定向、文件系统在这里就连成了一体。

\section{实验流程}

软件这边在 \lstinline|sys/software| 下 \lstinline|make| 即可，生成 \lstinline|main.coe| 和
\lstinline|main_d.coe| 两份存储器初始化文件，分别给指令存储器和数据存储器用。硬件这边在
\lstinline|sys/hardware| 下 \lstinline|make|，生成 \lstinline|top.bit|。把它们烧进板子后就能直接用。

这里稍微修改了提供的 \lstinline|Makefile|，没有使用课程提供的虚拟机，也没有自己编译一遍
RISC-V 工具链，而是直接用本机的 \lstinline|riscv64-unknown-elf-gcc|，再加上指定
架构的 \lstinline|-march=rv32i| 和 \lstinline|-mabi=ilp32|，感觉比用非常重的虚拟机，
或者自己处理编译的依赖要方便很多。

\section{实验总结与思考}

由于前面的几次实验已经完成了很大一部分内容，我也没有更换新的实现，所以这次实验的重点就只剩下
怎么把不同的硬件连在一起，再加上软件部分了。不过软件部分与正常的开发有很大区别，比如没有标准库，
所有的功能都需要自己写，还是有点麻烦的。

我在做硬件部分的时候没遇到什么问题，可能是最开始的框架搭的比较好，也没有多少需要改动的地方。
软件部分倒是遇到了一些问题，比如在最开始写的时候，我的代码经常没法编译成功，这也让我想到了
另一位同学遇到的问题：他没有把 \lstinline|entry| 函数放在 \lstinline|main.c| 的最顶部，
导致了一直出现的问题；改掉以后立刻就看起来正常了。

实际上我觉得挺费时的是调这个光标，可能是因为我在上个实验没有做这个。我最开始做的时候一直忘了
在各种地方进行 \lstinline|redraw|，导致显示出来的东西很诡异。要说我这个光标做的哪里不太好，
我只能说，首先我为了性能好，把原本的字模大小从 16x9 改成了 16x8，而光标我做的比较简陋，
显示的时候左右的间隙是不对称的，不太好看。其次，这里我用了 Home 和 End 键，这个本来是需要
专门处理长一点的控制码，但我偷懒了，直接把它们映射成了单个控制码，在现在看来没什么问题，
但是实际上是一个不太好的实现，如果有机会我倒是希望做一个完整的键盘出来。

另外的几个遗憾我也正好一并说出来吧。首先我没有做一个完整的文件系统，只是在内存上实现了，
当时想做的时候感觉接口真的太难写了，最后也没什么时间，就放弃了。第二个应该是需要在我的第一个
基础上实现，就是我希望做可运行的二进制文件。我感觉这个其实是很好做的，就运行的时候把 PC 移过去，
跑完了再回去就行，但是由于没有一个完整的文件系统，我做出来效果也不会好看，所以最后也没做。
（其实我感觉这个我应该做一下的）最后一个就是支持一下浮点数了，我尝试了这个，但是真的好难写，
硬件也要写很多，软件也要写很多，真的挺难的，最后我也没做出来。最开始想做的 GUI，
，由于我感觉要大改架构，也当然是没做出来了，只好无奈羡慕其他人。

所以最后也就完成了我一直想做的类似于 bash 的管道，和一个简单的文件系统来配合这个管道。
感觉还是时间不太够，端午节感觉对我影响很大，加上我的课比较多，最后也没有做成我预期的效果，
不过至少算是能看吧，完成了拓展内容，我也不能说做的很糟糕，至少还是有一点成就感的。

回顾整个实验的流程，我只能说，我似乎运气太好了一点，可能最开始所想的架构就不太容易出现问题，
在很多人碰到时序问题的时候我都没碰到，所以我在调试上花的时间大幅减少，主要集中在软件上，
做了一些我想做的东西，也算是顺利吧。后来我也帮其他同学调试了一下，我只能说，这个实验真的
好容易出问题，我能做成这样真的很幸运了。

我是和理论课一起选的这门，感觉两节课一起起到了相辅相成的作用，理论课上学的东西在这里能用上很多，
我个人感觉也比只使用 Logisim 更好理解理论课的很多内容。这门课真的很有意思，虽然后面的几个
实验都挺耗时间的，但做出来的成就感是远高于 Logisim 的。毕竟，谁不想做一个真实的计算机呢。

\section{致谢}

感谢李老师，解答了我的很多疑惑，在我遇到问题的时候也给了我很多帮助。
也感谢南大，给我们带来了这门让我觉得有意思并也有收获的课程。

感谢和我一起做实验的几个同学，大家都很热情，做实验的时候我心情很愉快。
感谢我的朋辈导师，他居然和我一起上课，这太巧了。
感谢我的一位学长，他的博客指引我把理论课和实践课一起选了，不然我不会出现在这里。

感谢雪碧和巧乐兹，做实验的时候经常吃这些，超级好吃。
感谢门尼和薛之谦，做实验的时候经常听他们，超级好听。

感谢 Vivado，感谢 Verilog，感谢 RISC-V，感谢所有的硬件工作者们。

感谢这个世界。

\end{document}
